\documentclass[12pt,a4paper]{article}

\usepackage[utf8]{inputenc}
\usepackage[T1]{fontenc}
\usepackage[english]{babel}
\usepackage[top=2.5cm,bottom=2.5cm,left=2.5cm,right=2.5cm]{geometry}
\usepackage{fancyhdr}
\setlength{\headheight}{14.5pt}
\pagestyle{fancy}
\fancyhf{}
\lhead{\small Università della Calabria}
\rhead{\small Data Warehouse and Visualization}
\cfoot{\thepage}
\renewcommand{\headrulewidth}{0.4pt}

\usepackage{booktabs}
\usepackage{array}
\usepackage{float}
\usepackage{caption}
\usepackage{xcolor}
\usepackage{listings}
\usepackage{amsmath}
\usepackage[colorlinks=true,linkcolor=blue,urlcolor=blue]{hyperref}
\setlength{\parskip}{6pt}
\setlength{\parindent}{0pt}

\lstset{
  basicstyle=\ttfamily\small,
  breaklines=true,
  backgroundcolor=\color{gray!8},
  frame=single,
  rulecolor=\color{gray!40},
  keywordstyle=\color{blue!70}\bfseries,
  commentstyle=\color{green!50!black},
  stringstyle=\color{red!60!black},
  showstringspaces=false,
  language=Python,
  xleftmargin=6pt,
  xrightmargin=6pt,
  aboveskip=8pt,
  belowskip=8pt
}

% ─────────────────────────────────────────────────────────────────────────────
\begin{document}
\thispagestyle{empty}

\begin{titlepage}
  \centering
  \vspace*{3cm}
  {\LARGE\bfseries Integrazione LLM nel Data Warehouse Olist\\[0.5em]
  Documentazione Tecnica degli Script \texttt{10\_script\_llm}}\\[1.5cm]
  {\large Emanuele Colecchia — Matricola 276527}\\[0.5cm]
  {\large Corso: Data Warehouse and Visualization}\\[0.5cm]
  {\large A.A. 2025/2026 — Università della Calabria}\\[2cm]
  {\normalsize\textit{Dataset: Brazilian E-Commerce Public Dataset (Olist, Kaggle)}}\\
  \vfill
  {\small Generato il \today}
\end{titlepage}

\tableofcontents
\newpage

% ═══════════════════════════════════════════════════════════════════════════
\section{Introduzione Teorica: LLM come Livello di Supporto}
% ═══════════════════════════════════════════════════════════════════════════

\subsection{Ruolo degli LLM nel processo DQ}

I Large Language Models (LLM) vengono utilizzati in questi script come
\textbf{livello di supporto} aggiuntivo rispetto alla pipeline classica di
Data Quality Assessment (DQA). Il principio fondamentale è che
\textit{gli LLM non sostituiscono la pipeline quantitativa}: gli scorecard DQ
rimangono la fonte autoritativa di verità, mentre gli LLM arricchiscono il
processo con:

\begin{itemize}
  \item \textbf{Interpretazione in linguaggio naturale} degli scorecard tecnici,
        traducendoli in report leggibili da stakeholder non tecnici.
  \item \textbf{Generazione di regole di validità domain-aware}: dato il contesto
        reale del dominio (e-commerce brasiliano, categorie prodotti, stati geografici),
        l'LLM suggerisce lambda Python specifiche per il dataset.
  \item \textbf{Classificazione della missingness}: supporto alla distinzione tra
        MCAR / MAR / MNAR con ragionamento domain-aware.
  \item \textbf{Pianificazione strutturata del cleaning}: generazione di piani JSON
        con priorità e metodi pandas specifici.
  \item \textbf{Narrazione degli audit log}: trasformazione di log tecnici in
        report comprensibili per il management.
\end{itemize}

\subsection{Provider LLM supportati}

Tutti e tre gli script supportano i seguenti provider, configurabili tramite
la variabile \texttt{LLM\_PROVIDER}:

\begin{table}[H]
  \centering
  \begin{tabular}{llll}
    \toprule
    \textbf{Provider} & \textbf{Modello} & \textbf{Requisito} & \textbf{Note} \\
    \midrule
    \texttt{openai}      & \texttt{gpt-4o-mini}              & API key OpenAI    & Cloud \\
    \texttt{anthropic}   & \texttt{claude-haiku-4-5-20251001} & API key Anthropic & Cloud \\
    \texttt{ollama}      & \texttt{llama3.2:3b}              & Server locale     & Default \\
    \texttt{huggingface} & \texttt{Qwen/Qwen2.5-1.5B-Instruct} & Token HF (opt.) & Cloud \\
    \bottomrule
  \end{tabular}
  \caption{Provider LLM supportati dalla pipeline}
\end{table}

\subsection{Architettura dei path}

Gli script sono stati adattati dalla versione originale (Google Colab +
Google Drive) a un'architettura \textbf{portable e locale}. I path vengono
calcolati automaticamente a partire dalla posizione dello script:

\begin{lstlisting}[language=Python]
_HERE   = os.path.dirname(os.path.abspath('__file__'))  # consegna/10_script_llm/
_ROOT   = os.path.dirname(_HERE)                         # consegna/
RAW_PATH     = os.path.join(_ROOT, '1_raw_data')    + os.sep
CLEANED_PATH = os.path.join(_ROOT, '3_cleaned_data') + os.sep
OUTPUT_PATH  = os.path.join(_ROOT, '10_script_llm', 'outputs') + os.sep
\end{lstlisting}

Questo garantisce che gli script funzionino su qualsiasi macchina senza
modifiche manuali ai path.

% ═══════════════════════════════════════════════════════════════════════════
\section{Modifiche Apportate — Adattamento al Dataset Olist}
% ═══════════════════════════════════════════════════════════════════════════

\subsection{Problemi nella versione originale}

La versione originale dei notebook presentava i seguenti problemi che
impedivano l'esecuzione:

\begin{enumerate}
  \item \textbf{Path Google Colab non portabili:}
        \texttt{BASE\_PATH = '/content/drive/Shareddrives/...'} e
        \texttt{drive.mount('/content/drive')} — inutilizzabili al di fuori di Colab.
  \item \textbf{Dataset generici non correlati a Olist:}
        \begin{itemize}
          \item L1: \texttt{orders\_table.csv} (dati sintetici generici)
          \item L2: \texttt{patient\_readings.csv} (dati medici, nessuna relazione con il progetto)
          \item L3: \texttt{L3\_products\_dirty.csv} (dati retail generici)
        \end{itemize}
  \item \textbf{Prompt LLM con contesto sbagliato:} riferimenti a dataset medici,
        colonne inesistenti (\texttt{admission\_date}, \texttt{category}, \texttt{brand}).
  \item \textbf{Celle di installazione sistema inutili:} \texttt{!apt-get install zstd},
        \texttt{!curl -fsSL https://ollama.com/install.sh | sh} — specifiche di Colab.
  \item \textbf{Codice duplicato:} funzione \texttt{call\_llm\_ollama} definita più volte,
        poi ridefinita come \texttt{call\_llm} — confusione e ridondanza.
\end{enumerate}

\subsection{Interventi effettuati}

\begin{table}[H]
  \centering
  \begin{tabular}{p{4cm}p{5cm}p{5cm}}
    \toprule
    \textbf{Problema} & \textbf{Versione originale} & \textbf{Versione adattata} \\
    \midrule
    Path & \texttt{/content/drive/Shareddrives/...} & Path relativi con \texttt{os.path} \\
    Dataset L1 & \texttt{orders\_table.csv} (sintetico) & \texttt{olist\_orders\_dataset.csv} \\
    Dataset L2 & \texttt{patient\_readings.csv} (medico) & \texttt{olist\_orders\_dataset.csv} + \texttt{olist\_order\_items\_dataset.csv} \\
    Dataset L3 & \texttt{L3\_products\_dirty.csv} (generico) & \texttt{olist\_products\_dataset.csv} + \texttt{dim\_products.csv} \\
    Scorecard & File generici & \texttt{4\_dq\_scorecards/dq\_scorecard\_*.csv} reali \\
    Colab setup & \texttt{drive.mount()}, \texttt{!apt-get} & Rimossi \\
    Prompt LLM & Contesto generico/medico & Contesto Olist (BRL, 2016-2018, BR) \\
    Funzione LLM & Duplicata e incompleta & Unificata in \texttt{call\_llm()} completa \\
    \bottomrule
  \end{tabular}
  \caption{Modifiche apportate per adattare i notebook al dataset Olist}
\end{table}

% ═══════════════════════════════════════════════════════════════════════════
\section{L1 — LLM-Supported Data Quality Assessment}
% ═══════════════════════════════════════════════════════════════════════════

\subsection{Descrizione}

Lo script L1 utilizza gli LLM per analizzare gli scorecard DQ già prodotti
dalla pipeline classica (\texttt{data\_quality\_assessment\_LOCAL.ipynb}) e
produrre tre output in linguaggio naturale:

\begin{enumerate}
  \item \textbf{Audit Report}: interpretazione professionale dello scorecard per stakeholder.
  \item \textbf{Regole di validità}: generazione di lambda Python domain-aware.
  \item \textbf{Sommario stakeholder}: spiegazione semplificata per management non tecnico.
\end{enumerate}

\subsection{Input}

\begin{itemize}
  \item \texttt{1\_raw\_data/olist\_orders\_dataset.csv} (99.441 righe)
  \item \texttt{1\_raw\_data/olist\_customers\_dataset.csv} (99.441 righe)
  \item \texttt{1\_raw\_data/olist\_order\_items\_dataset.csv} (112.650 righe)
  \item \texttt{1\_raw\_data/olist\_products\_dataset.csv} (32.951 righe)
  \item \texttt{1\_raw\_data/olist\_sellers\_dataset.csv} (3.095 righe)
  \item \texttt{1\_raw\_data/olist\_order\_payments\_dataset.csv} (103.886 righe)
  \item \texttt{1\_raw\_data/olist\_order\_reviews\_dataset.csv} (99.224 righe)
  \item \texttt{1\_raw\_data/olist\_geolocation\_dataset.csv} (1.000.163 righe)
  \item \texttt{4\_dq\_scorecards/dq\_scorecard\_*.csv} (scorecard per ogni tabella)
\end{itemize}

\subsection{Output}

\begin{itemize}
  \item \texttt{10\_script\_llm/outputs/dq\_audit\_\{table\}.md} — report audit per ogni tabella
  \item \texttt{10\_script\_llm/outputs/dq\_generated\_rules\_\{table\}.md} — regole validità
  \item \texttt{10\_script\_llm/outputs/stakeholder\_summaries.md} — sommari stakeholder
\end{itemize}

\subsection{Prompt LLM — Use Case 1: Audit Report}

\begin{lstlisting}[language=Python, caption=System prompt per audit report]
SYSTEM_AUDIT = """You are a senior Data Quality Analyst specializing in Data Warehouse
systems for large-scale e-commerce platforms. The dataset is from Olist, the largest
Brazilian online marketplace. Your task is to interpret a structured DQ scorecard and
produce a clear, actionable audit report for a business stakeholder audience.
Be specific about root causes and prioritize recommendations by business impact.
Write in Italian."""
\end{lstlisting}

\begin{lstlisting}[language=Python, caption=User prompt per audit report (template)]
PROMPT_AUDIT = f"""Di seguito e' riportato lo scorecard DQ per la tabella '{table_name}'
del Data Warehouse Olist (e-commerce brasiliano, ~100.000 ordini reali, 2016-2018).
Analizza e produci un report strutturato con:
1. Riepilogo Esecutivo (2-3 frasi)
2. Problemi Critici (richiedono intervento immediato)
3. Problemi Moderati (da risolvere prima del prossimo ciclo ETL)
4. Priorita' di Cleaning Raccomandata (lista ordinata per impatto)
5. Rischio Stimato per Analytics Downstream

SCORECARD:
{scorecard_json}
"""
\end{lstlisting}

\subsection{Prompt LLM — Use Case 2: Generazione Regole di Validità}

\begin{lstlisting}[language=Python, caption=System prompt per generazione regole]
SYSTEM_RULES = """You are a Data Engineer expert in Data Quality for e-commerce DW systems.
Given column metadata and sample values from the Olist Brazilian e-commerce platform,
generate Python pandas validation rule functions (lambdas) using only standard
pandas/numpy operations.
Each rule must be a lambda that takes a Series and returns a boolean Series
(True = valid).
Return ONLY valid Python code in a dict named 'validity_rules'.
Use domain knowledge: Brazilian state codes (2 uppercase letters), UUID-format IDs,
positive prices, order_status in a defined set, etc."""
\end{lstlisting}

\begin{lstlisting}[language=Python, caption=Esempio output atteso dalle regole generate]
validity_rules = {
    'order_id':      lambda s: s.str.len() == 32,
    'customer_id':   lambda s: s.str.len() == 32,
    'order_status':  lambda s: s.isin(['delivered','shipped','processing',
                                        'canceled','unavailable','invoiced',
                                        'created','approved']),
    'order_purchase_timestamp': lambda s: pd.to_datetime(s, errors='coerce').notna(),
}
\end{lstlisting}

% ═══════════════════════════════════════════════════════════════════════════
\section{L2 — LLM-Supported Missing Values \& Outlier Analysis}
% ═══════════════════════════════════════════════════════════════════════════

\subsection{Descrizione}

Lo script L2 analizza i pattern di valori mancanti e gli outlier nei dataset
Olist, con supporto LLM per:

\begin{enumerate}
  \item \textbf{Interpretazione MCAR/MAR/MNAR}: classificazione domain-aware
        della tipologia di missingness.
  \item \textbf{Strategia di imputazione}: generazione di codice Python per
        la gestione ottimale dei NULL.
  \item \textbf{Narrativa outlier}: spiegazione business degli outlier rilevati.
\end{enumerate}

\subsection{Dataset analizzati}

\begin{itemize}
  \item \texttt{olist\_orders\_dataset.csv}: missingness su \texttt{order\_approved\_at}
        (160 NULL), \texttt{order\_delivered\_carrier\_date} (1.783 NULL),
        \texttt{order\_delivered\_customer\_date} (2.965 NULL).
  \item \texttt{olist\_order\_items\_dataset.csv}: outlier su \texttt{price}
        (max 6.735 BRL) e \texttt{freight\_value}.
  \item \texttt{olist\_products\_dataset.csv}: outlier su misure fisiche
        (es. peso 40.000g), missingness su colonne dimensionali.
  \item \texttt{olist\_customers\_dataset.csv}: analisi completezza e unicità.
  \item \texttt{olist\_sellers\_dataset.csv}: analisi distribuzione geografica e outlier.
  \item \texttt{olist\_order\_payments\_dataset.csv}: outlier su \texttt{payment\_value}
        e \texttt{payment\_installments}.
  \item \texttt{olist\_order\_reviews\_dataset.csv}: missingness su
        \texttt{review\_comment\_title} e \texttt{review\_comment\_message}.
  \item \texttt{olist\_geolocation\_dataset.csv}: outlier su coordinate lat/lng.
\end{itemize}

\subsection{Classificazione Missingness}

La classe \texttt{MissingnessAnalyzer} applica un'euristica basata su:

\begin{itemize}
  \item \textbf{MCAR}: percentuale di missingness $< 0.5\%$ — probabile casualità.
  \item \textbf{MAR}: missingness correlata a variabili categoriche osservabili
        (variazione $> 5\%$ del tasso di missingness tra categorie).
  \item \textbf{MNAR}: nessuna correlazione rilevata — possibile null semantico
        (es. ordine non ancora consegnato).
\end{itemize}

Per il dataset Olist, i tre campi mancanti sono classificati come \textbf{MNAR
(null semantici)}: la mancanza di \texttt{order\_delivered\_customer\_date} non
è un errore — significa semplicemente che l'ordine non è ancora stato consegnato.

\subsection{Prompt LLM — Interpretazione Missingness}

\begin{lstlisting}[language=Python, caption=System prompt per interpretazione missingness]
SYSTEM_MISSING = """You are a senior Data Engineer specializing in e-commerce
Data Warehouses. You are analyzing the Olist dataset, a Brazilian e-commerce
platform with ~100,000 orders from 2016-2018.
Your role is to interpret missing value patterns, classify them as MCAR/MAR/MNAR
with domain knowledge, and suggest appropriate imputation or handling strategies.
Write in Italian. Be specific and use domain context (e.g., orders without
delivery dates may simply not have been delivered yet -- semantic nulls,
not data errors)."""
\end{lstlisting}

\begin{lstlisting}[language=Python, caption=User prompt per interpretazione missingness (estratto)]
PROMPT_MISSING = f"""Tabella: olist_orders (99.441 ordini e-commerce brasiliani, 2016-2018)
Colonne con valori mancanti:
{miss_json}

Nota di dominio:
- order_approved_at: data di approvazione pagamento. NULL se il pagamento
  non e' stato ancora approvato.
- order_delivered_carrier_date: data di consegna al corriere. NULL se
  l'ordine e' ancora in lavorazione.
- order_delivered_customer_date: data di consegna al cliente. NULL se
  non ancora consegnato.

Per ogni colonna:
1. Classifica il tipo di missingness (MCAR/MAR/MNAR) con motivazione
2. Spiega l'impatto sul DW se i NULL non vengono gestiti correttamente
3. Raccomanda la strategia ottimale (flag binario, esclusione, imputazione)
"""
\end{lstlisting}

\subsection{Outlier Detection — Metodi}

\begin{table}[H]
  \centering
  \begin{tabular}{lll}
    \toprule
    \textbf{Metodo} & \textbf{Formula} & \textbf{Uso} \\
    \midrule
    IQR ($k=1.5$) & $[Q_1 - 1.5 \cdot IQR,\; Q_3 + 1.5 \cdot IQR]$ & Default \\
    Z-score ($\sigma=3$) & $|z| > 3$ dove $z = (x - \mu)/\sigma$ & Distribuz.\ normali \\
    \bottomrule
  \end{tabular}
  \caption{Metodi di outlier detection implementati in OutlierDetector}
\end{table}

\subsection{Prompt LLM — Narrativa Outlier}

\begin{lstlisting}[language=Python, caption=System prompt per narrativa outlier]
SYSTEM_OUTLIER = """You are a senior Data Analyst for Olist, a Brazilian
e-commerce marketplace. You interpret outlier detection results and provide
business-context narrative explanations. For each detected outlier:
- Explain potential real-world causes (luxury products, errors, fraud, etc.)
- Assess business impact (distorted averages, wrong recommendations, etc.)
- Recommend the appropriate handling strategy
Write in Italian."""
\end{lstlisting}

% ═══════════════════════════════════════════════════════════════════════════
\section{L3 — LLM-Supported Data Cleaning}
% ═══════════════════════════════════════════════════════════════════════════

\subsection{Descrizione}

Lo script L3 utilizza gli LLM per supportare il processo di cleaning di
tutte le 8 tabelle Olist, con due use case principali applicati sistematicamente:

\begin{enumerate}
  \item \textbf{Piano di cleaning strutturato}: l'LLM genera per ogni tabella
        un piano JSON con step numerati, metodo pandas da applicare, priorità
        (high/medium/low) e impatto sul business se lo step non viene eseguito.
  \item \textbf{Narrativa audit}: dato il confronto statistico prima/dopo
        cleaning, l'LLM produce una narrativa professionale delle trasformazioni
        effettuate, quantificando il miglioramento della qualità.
\end{enumerate}

Per la tabella \texttt{products}, è presente anche un terzo use case specifico:
mapping delle categorie dal portoghese all'inglese, con verifica del mapping
ufficiale Olist e integrazione LLM per le categorie non coperte.

\subsection{Input/Output}

\begin{table}[H]
  \centering
  \begin{tabular}{lll}
    \toprule
    \textbf{File} & \textbf{Tipo} & \textbf{Contenuto} \\
    \midrule
    \texttt{olist\_products\_dataset.csv} & Input raw & 32.951 prodotti sporchi \\
    \texttt{product\_category\_name\_translation.csv} & Reference & 71 mapping PT$\to$EN \\
    \texttt{dim\_products.csv} & Input cleaned & Versione pulita post-pipeline \\
    \texttt{audit\_log.csv} & Input & Log trasformazioni \\
    \texttt{L3\_category\_mapping.json} & Output & Mapping finale PT$\to$EN \\
    \texttt{L3\_cleaning\_plan.json} & Output & Piano cleaning strutturato \\
    \texttt{L3\_audit\_narrative.md} & Output & Narrativa audit in italiano \\
    \bottomrule
  \end{tabular}
  \caption{File di input e output di L3}
\end{table}

\subsection{Prompt LLM — Use Case 1: Mapping Categorie}

\begin{lstlisting}[language=Python, caption=System prompt per mapping categorie PT→EN]
SYSTEM_ENUM = """You are a Data Engineer building a product taxonomy for a
Brazilian e-commerce DW. The platform is Olist, an online marketplace
aggregating multiple sellers.
You receive Portuguese product category names and must map them to standardized
English equivalents.
Rules:
- Use standard e-commerce taxonomy (Electronics, Home & Garden, Sports,
  Beauty & Health, etc.)
- Return ONLY a Python dict named 'category_mapping' with 'pt_name':'en_name'
- If a category is 'nan' or invalid, map it to 'unknown'
- Be consistent: same concept -> same English label"""
\end{lstlisting}

Il mapping ufficiale contiene 71 coppie PT$\to$EN. Il sistema usa
\texttt{product\_category\_name\_translation.csv} come ground truth e invoca
l'LLM solo per categorie eventualmente non mappate nel file ufficiale.

\subsection{Prompt LLM — Use Case 2: Piano di Cleaning}

\begin{lstlisting}[language=Python, caption=System prompt per piano cleaning]
SYSTEM_PLAN = """You are a senior Data Engineer building ETL cleaning pipelines
for a Brazilian e-commerce Data Warehouse (Olist). You produce structured,
actionable cleaning plans in JSON format that can be directly translated into
Python/pandas code. Prioritize by business impact. Be specific about the
pandas method to use. Return ONLY valid JSON."""
\end{lstlisting}

\begin{lstlisting}[language=Python, caption=Struttura JSON attesa dal piano cleaning]
{
  "table": "olist_products_dataset",
  "steps": [
    {
      "step_id": 1,
      "name": "fill_missing_category",
      "description": "Imputa le categorie mancanti con 'unknown'",
      "columns": ["product_category_name"],
      "pandas_method": "fillna('unknown')",
      "priority": "high",
      "business_impact": "Categorie NULL rendono impossibile l'analisi fatturato per categoria"
    },
    {
      "step_id": 2,
      "name": "add_english_category",
      "description": "Aggiunge colonna product_category_name_english tramite JOIN",
      "columns": ["product_category_name"],
      "pandas_method": "map(cat_translation_dict)",
      "priority": "high",
      "business_impact": "Dashboard Tableau richiede categorie in inglese"
    }
  ]
}
\end{lstlisting}

% ═══════════════════════════════════════════════════════════════════════════
\section{Istruzioni per l'Esecuzione}
% ═══════════════════════════════════════════════════════════════════════════

\subsection{Prerequisiti}

\begin{lstlisting}[language=bash, caption=Installazione dipendenze Python]
pip install openai anthropic scikit-learn scipy requests --quiet
\end{lstlisting}

Per usare \textbf{Ollama} (provider consigliato, locale, gratuito):

\begin{lstlisting}[language=bash, caption=Installazione e avvio Ollama]
# 1. Installa Ollama (Mac/Linux)
curl -fsSL https://ollama.com/install.sh | sh

# 2. Scarica il modello
ollama pull llama3.2:3b

# 3. Avvia il server (si avvia automaticamente alle richieste API)
ollama serve
\end{lstlisting}

\subsection{Ordine di esecuzione}

\begin{enumerate}
  \item \textbf{L1}: apri \texttt{L1\_LLM\_data\_quality\_assessment.ipynb} e
        esegui tutte le celle. Richiede che gli scorecard DQ siano già in
        \texttt{4\_dq\_scorecards/} (generati da \texttt{data\_quality\_assessment\_LOCAL.ipynb}).
  \item \textbf{L2}: apri \texttt{L2\_Colab2\_LLM\_Missing\_Outliers.ipynb} e
        esegui tutte le celle. Legge direttamente i CSV raw da \texttt{1\_raw\_data/}.
  \item \textbf{L3}: apri \texttt{L3\_Colab2\_LLM\_Cleaning.ipynb} e
        esegui tutte le celle. Richiede sia i CSV raw che i cleaned da
        \texttt{3\_cleaned\_data/}.
\end{enumerate}

\subsection{Output generati}

Tutti gli output vengono salvati in \texttt{10\_script\_llm/outputs/}:

\begin{table}[H]
  \centering
  \begin{tabular}{lll}
    \toprule
    \textbf{Script} & \textbf{File} & \textbf{Contenuto} \\
    \midrule
    L1 & \texttt{dq\_audit\_\{table\}.md} & Audit report per ogni tabella \\
    L1 & \texttt{dq\_generated\_rules\_\{table\}.md} & Regole di validità Python \\
    L1 & \texttt{stakeholder\_summaries.md} & Sommari per management \\
    L2 & \texttt{L2\_missing\_interpretation.md} & Interpretazione MCAR/MAR/MNAR \\
    L2 & \texttt{L2\_outlier\_narrative.md} & Narrativa outlier \\
    L2 & \texttt{L2\_imputation\_code.py} & Codice imputazione generato \\
    L3 & \texttt{L3\_category\_mapping.json} & Mapping PT$\to$EN finale \\
    L3 & \texttt{L3\_cleaning\_plan.json} & Piano cleaning strutturato \\
    L3 & \texttt{L3\_audit\_narrative.md} & Narrativa audit prodotti \\
    \bottomrule
  \end{tabular}
  \caption{Output generati dalla pipeline LLM}
\end{table}

% ═══════════════════════════════════════════════════════════════════════════
\section{Considerazioni Metodologiche}
% ═══════════════════════════════════════════════════════════════════════════

\subsection{Limiti degli LLM in questo contesto}

\begin{itemize}
  \item \textbf{Non determinismo}: la stessa chiamata LLM può produrre output
        diversi tra esecuzioni successive. Per questo il codice generato viene
        sempre validato prima dell'esecuzione (\texttt{exec()} in sandbox).
  \item \textbf{Allucinazioni}: l'LLM può generare regole di validità plausibili
        ma errate (es.\ formato errato per \texttt{order\_id}). Le regole generate
        devono essere revisionate manualmente prima di essere usate in produzione.
  \item \textbf{Dipendenza dal modello}: modelli più piccoli (l  \item \textbf{Dipendenza dal modello}: modelli più piccoli (llama3.2:3b) producono
        output meno precisi rispetto a GPT-4o o Claude Opus. Per produzione è
        consigliato usare un modello cloud.
\end{itemize}

\subsection{Sicurezza del codice generato}

Tutti i blocchi di codice Python generati dall'LLM vengono eseguiti in un
ambiente sandbox isolato:

\begin{lstlisting}[language=Python, caption=Esecuzione sicura del codice LLM]
local_env = {}   # ambiente isolato -- non accede a variabili globali
try:
    exec(code_str, {}, local_env)
    result = local_env.get('category_mapping', {})
except Exception as e:
    print(f'Parsing LLM code fallito: {e}')
    result = {}
\end{lstlisting}

\subsection{Conclusioni}

Gli script L1, L2 e L3 dimostrano un approccio ibrido alla Data Quality:
la pipeline quantitativa classica fornisce numeri oggettivi e riproducibili,
mentre gli LLM aggiungono interpretazione domain-aware, generazione di codice
e comunicazione verso stakeholder non tecnici. L'adattamento al dataset Olist
rende l'intera pipeline eseguibile localmente senza dipendenze da Google Colab
o da API chiave per il provider Ollama.

\end{document}
